\subsection{Gestione del carrello}

Il \hyperlink{glossario:carrello}{carrello} \`e la sezione in cui puoi vedere tutti i prodotti che hai aggiunto durante l'utilizzo della chat.\\
Si trova nella parte destra della schermata e viene aggiornato automaticamente ogni volta che fai una richiesta, ad esempio aggiungendo o modificando un prodotto.

\begin{figure}[H]
	\centering
	\includegraphics[height=0.5\textheight,keepaspectratio]{ManualeUtente/carrello}
	\caption{Area carrello}
\end{figure}

Puoi anche espandere o nascondere il \hyperlink{glossario:carrello}{carrello} utilizzando il pulsante con l'icona del \hyperlink{glossario:carrello}{carrello}, cos\`i da avere pi\`u spazio per lavorare nella \hyperlink{glossario:chat}{chat} quando serve.

\subsubsection{Visualizzazione prodotti}

All'interno del \hyperlink{glossario:carrello}{carrello} trovi l'elenco completo dei prodotti selezionati.\\
Per ogni prodotto vengono mostrate le informazioni principali:
\begin{itemize}[itemsep=5pt, parsep=5pt, label=$\scriptstyle\bullet$]
	\item Nome del prodotto
	\item Codice identificativo
	\item Quantit\`a
	\item Prezzo, unitario e totale
\end{itemize}

In questo modo puoi controllare facilmente cosa hai inserito nel tuo ordine.

\subsubsection{Aggiornamento del carrello}

Il \hyperlink{glossario:carrello}{carrello} si aggiorna automaticamente in base a ci\`o che scrivi nella chat.\\
Ad esempio, puoi:
\begin{itemize}[itemsep=5pt, parsep=5pt, label=$\scriptstyle\bullet$]
	\item Aggiungere un prodotto semplicemente chiedendolo nella \hyperlink{glossario:chat}{chat}
	\item Modificare la quantit\`a di un prodotto gi\`a presente
\end{itemize}

Il sistema interpreta la tua richiesta e aggiorna subito il carrello.

Non \`e quindi necessario interagire direttamente con il \hyperlink{glossario:carrello}{carrello}: tutte le operazioni vengono gestite tramite la \hyperlink{glossario:chat}{chat}, in modo semplice e veloce.

\subsubsection{Totale ordine}

Nella parte inferiore del \hyperlink{glossario:carrello}{carrello} viene mostrato il totale dell'ordine.

Il valore si aggiorna automaticamente ogni volta che modifichi il contenuto del \hyperlink{glossario:carrello}{carrello}, cos\`i hai sempre sotto controllo il costo complessivo.

\subsubsection{Invio ordine}

Quando hai completato il tuo \hyperlink{glossario:ordine}{ordine}, puoi inviarlo cliccando sul pulsante \vr{Invia \hyperlink{glossario:ordine}{ordine}}.

A questo punto viene mostrata una finestra di conferma, che riassume le informazioni principali dell'\hyperlink{glossario:ordine}{ordine}, come il numero di prodotti e il totale.

\begin{figure}[H]
	\centering
	\includegraphics[width=0.5\textwidth]{ManualeUtente/conferma_ordine}
	\caption{Finestra di conferma ordine}
\end{figure}

Puoi quindi:
\begin{itemize}[itemsep=5pt, parsep=5pt, label=$\scriptstyle\bullet$]
	\item Confermare l'\hyperlink{glossario:ordine}{ordine}, per completare l'operazione
	\item Annullare, per tornare al \hyperlink{glossario:carrello}{carrello} e apportare eventuali modifiche
\end{itemize}

Questa conferma ti permette di controllare tutto prima di inviare definitivamente l'ordine.
